%==========================================================================
% BCC502 -- Metodologia Científica para Ciência da Computação
% Template do exercício integrador: o seu TCC em seis slides
%
% Regras:
%  1. Preencha apenas \titulotrabalho, \nomealuno, \matricula e o conteúdo entre colchetes.
%  2. Mantenha a capa e os seis slides, na ordem dada. Não crie slide de
%     sumário nem de referências nem outros slides.
%  3. Não altere os títulos dos slides, as cores nem o rodapé.
%  4. Apague todo texto entre colchetes vermelhos antes de entregar.
%     No slide 2, use de 3 a 5 trabalhos e de 3 a 5 características.
%  5. Entregue um único PDF, nomeado matricula_nome.pdf.
%==========================================================================
\documentclass[aspectratio=169,11pt]{beamer}

\usepackage[utf8]{inputenc}
\usepackage[T1]{fontenc}
\usepackage[brazil]{babel}
\usepackage{microtype}
\usepackage{csquotes}
\usepackage{booktabs}
\usepackage{array}
\usepackage{tabularx}
\usepackage[most]{tcolorbox}

%--------------------------------------------------------------------------
% DADOS DO ALUNO (preencha)
%--------------------------------------------------------------------------
\newcommand{\titulotrabalho}{Título provisório do trabalho}
\newcommand{\nomealuno}{Nome Completo do Aluno}
\newcommand{\matricula}{00.0.0000}

%--------------------------------------------------------------------------
% Cores e tema (não alterar)
%--------------------------------------------------------------------------
\definecolor{mtext}{RGB}{38, 38, 38}
\definecolor{mgray}{RGB}{120, 120, 120}
\definecolor{mlight}{RGB}{215, 215, 215}
\definecolor{maccent}{HTML}{5D5179}
\definecolor{mplace}{RGB}{176, 58, 46}

\usetheme{default}
\usecolortheme{default}
\setbeamercolor{normal text}{fg=mtext, bg=white}
\setbeamercolor{frametitle}{fg=mtext}
\setbeamercolor{itemize item}{fg=maccent}
\setbeamercolor{itemize subitem}{fg=mgray}

\setbeamerfont{frametitle}{series=\bfseries, size=\large}
\setbeamerfont{captitulo}{series=\bfseries, size=\fontsize{22}{27}}

\setbeamersize{text margin left=1.3cm, text margin right=1.3cm}
\setbeamertemplate{navigation symbols}{}
\setbeamertemplate{itemize items}[circle]
\setlength{\parskip}{0pt}

% Título: texto à esquerda e um filete curto de destaque
\setbeamertemplate{frametitle}{%
  \vskip 0.15cm
  {\usebeamerfont{frametitle}\usebeamercolor[fg]{frametitle}\insertframetitle\par}
  \vskip 0.1cm
  \textcolor{maccent}{\rule{1.1cm}{1.5pt}}
  \vskip -0.2cm
}

% Rodapé discreto
\setbeamertemplate{footline}{%
  \leavevmode%
  \hbox{%
    \begin{beamercolorbox}[wd=0.8\paperwidth, ht=2.5ex, dp=1.2ex, leftskip=1.3cm]{normal text}%
      \scriptsize\color{mgray} BCC502 \textperiodcentered{} \nomealuno{} \textperiodcentered{} Matrícula \matricula
    \end{beamercolorbox}%
    \begin{beamercolorbox}[wd=0.2\paperwidth, ht=2.5ex, dp=1.2ex, rightskip=1.3cm, right]{normal text}%
      \scriptsize\color{mgray}\insertframenumber{} / 6
    \end{beamercolorbox}%
  }%
}

% Bloco sem preenchimento, só com traço lateral
\newtcolorbox{bloco}[1]{%
  enhanced, frame hidden, boxrule=0pt, arc=0pt,
  colback=white, colbacktitle=white,
  borderline west={1.5pt}{0pt}{maccent},
  left=7pt, right=2pt, top=3pt, bottom=3pt,
  before skip=5pt, after skip=2pt,
  before upper={\textcolor{maccent}{\small\bfseries #1}\par\vspace{0.2em}},
}

% Rótulo de coluna ou de lista
\newcommand{\rotulo}[1]{\textcolor{maccent}{\small\bfseries #1}}

% Texto a ser substituído pelo aluno
\newcommand{\preencher}[1]{\textcolor{mplace}{\textit{[#1]}}}

% Tabelas: apenas linhas horizontais finas
\setlength{\heavyrulewidth}{0.8pt}
\setlength{\lightrulewidth}{0.4pt}
\newcolumntype{L}{>{\raggedright\arraybackslash}X}

\begin{document}

%==========================================================================
% CAPA (não conta como slide de etapa)
%==========================================================================
\begin{frame}[plain,c]
  \relax
  {\usebeamercolor[fg]{frametitle}\usebeamerfont{captitulo}\titulotrabalho\par}
  \vskip 0.35cm
  \textcolor{maccent}{\rule{1.1cm}{1.5pt}}
  \vskip 0.5cm
  {\large \nomealuno\par}
  \vskip 0.15cm
  {\small\color{mgray} BCC502 \textperiodcentered{} Matrícula \matricula\par}
\end{frame}
\setcounter{framenumber}{0}

%==========================================================================
% SLIDE 1
%==========================================================================
\begin{frame}[t]{1. Ideia crua}
  \begin{bloco}{O que pretendo fazer (uma ou duas frases, sem polir)}
    \preencher{Escreva a ideia como ela está hoje, mesmo que vaga.}
  \end{bloco}

  \vspace{0.5cm}
  \begin{columns}[T]
    \column{0.31\textwidth}
      \rotulo{Termos avaliativos sem métrica}
      \begin{itemize}\small
        \item \preencher{ex.: \enquote{melhor}}
        \item \preencher{\dots}
      \end{itemize}
    \column{0.31\textwidth}
      \rotulo{Conceitos sem definição}
      \begin{itemize}\small
        \item \preencher{conceito}
        \item \preencher{\dots}
      \end{itemize}
    \column{0.31\textwidth}
      \rotulo{Comparadores com subtipos não examinados}
      \begin{itemize}\small
        \item \preencher{comparador}
        \item \preencher{\dots}
      \end{itemize}
  \end{columns}
\end{frame}

%==========================================================================
% SLIDE 2
%==========================================================================
\begin{frame}[t]{2. Mapa da área}
  \footnotesize
  \renewcommand{\arraystretch}{1.35}
  \begin{tabularx}{\textwidth}{@{}>{\raggedright\arraybackslash}p{0.27\textwidth} *{4}{>{\centering\arraybackslash}X}@{}}
    \toprule
    \textcolor{mgray}{Trabalho ou artefato} & \preencher{Caract.\ 1} & \preencher{Caract.\ 2} & \preencher{Caract.\ 3} & \preencher{Caract.\ 4} \\
    \midrule
    \preencher{Artefato A (referência)} & X & X &   & ? \\
    \preencher{Artefato B (referência)} & X &   &   & X \\
    \preencher{Artefato C (referência)} &   & X & X & ? \\
    \bottomrule
  \end{tabularx}

  \vspace{0.5cm}
  \begin{itemize}\small
    \item \rotulo{Como escolhi os artefatos}\quad \preencher{critério objetivo, ex.: citados em pelo menos dois artigos}
    \item \rotulo{Por que as características são relevantes}\quad \preencher{que dificuldade dos trabalhos atuais cada uma responde}
    \item \rotulo{?}\quad informação que não consegui verificar
  \end{itemize}
\end{frame}

%==========================================================================
% SLIDE 3
%==========================================================================
\begin{frame}[t]{3. Degrau e teto da afirmação}
  \begin{bloco}{Degrau que o meu TCC pode alcançar}
    \preencher{1 a 5, com o nome do degrau (tabela do handout)}
  \end{bloco}

  \begin{bloco}{A afirmação mais forte que consigo sustentar em um semestre}
    \preencher{Uma frase.}
  \end{bloco}

  \begin{bloco}{Por que não o degrau acima}
    \begin{itemize}\small
      \item \textbf{Bases de dados e métricas consensuais:} \preencher{existem ou faltam}
      \item \textbf{Trabalhos de referência recentes:} \preencher{existem ou faltam}
      \item \textbf{Tempo e recursos:} \preencher{o que limita}
    \end{itemize}
  \end{bloco}
\end{frame}

%==========================================================================
% SLIDE 4
%==========================================================================
\begin{frame}[t]{4. Pergunta}
  \begin{columns}[T]
    \column{0.55\textwidth}
      \footnotesize
      \renewcommand{\arraystretch}{1.2}
      \begin{tabularx}{\textwidth}{@{}l L@{}}
        \toprule
        \textcolor{mgray}{Elemento} & \textcolor{mgray}{No meu TCC} \\
        \midrule
        População & \preencher{\dots} \\
        Intervenção & \preencher{\dots} \\
        Comparação (controle) & \preencher{\dots e por que este controle} \\
        Resultados & \preencher{métricas definidas} \\
        Contexto & \preencher{\dots} \\
        Design do experimento & \preencher{\dots} \\
        \bottomrule
      \end{tabularx}
    \column{0.40\textwidth}
      \footnotesize
      \begin{itemize}
        \item \rotulo{Família e subtipo:} \preencher{\dots}
        \item \rotulo{Perguntas anteriores necessárias:} \preencher{\dots}
        \item \rotulo{O mapa do slide 2 responde alguma?} \preencher{sim ou não, qual}
        \item \rotulo{Compatível com o degrau do slide 3?} \preencher{sim ou não. Se não, o que ajustei}
      \end{itemize}
  \end{columns}

  \vspace{0cm}
  \begin{bloco}{Pergunta de pesquisa}
    \preencher{A pergunta, escrita por extenso.}
  \end{bloco}
\end{frame}

%==========================================================================
% SLIDE 5
%==========================================================================
\begin{frame}[t]{5. Hipótese e refutação}
  \begin{bloco}{Hipótese}
    \preencher{Uma frase: o que acredito que seja verdade e ainda não está demonstrado.}
  \end{bloco}

  \begin{bloco}{Resultado que a refutaria}
    \preencher{Que dado ou resultado me obrigaria a abandonar a hipótese.}
  \end{bloco}

  \begin{bloco}{O que faço se ela for refutada}
    \preencher{Plano alternativo e o que ainda seria contribuição.}
  \end{bloco}
\end{frame}

%==========================================================================
% SLIDE 6
%==========================================================================
\begin{frame}[t]{6. Antes e depois}
  \begin{columns}[T]
    \column{0.47\textwidth}
      \begin{bloco}{Antes (slide 1)}
        \small\preencher{Ideia crua.}
      \end{bloco}
    \column{0.47\textwidth}
      \begin{bloco}{Depois (slide 4)}
        \small\preencher{Pergunta final.}
      \end{bloco}
  \end{columns}

  \vspace{0.3cm}
  \begin{bloco}{O que mudou}
    \small
    \preencher{Um parágrafo curto: o que abandonei ou rebaixei, e o que mudou.}
  \end{bloco}
\end{frame}

\end{document}
